% ============================================================
% ANHANG BCLM‑LL: EXPERIMENTELLES PROTOKOLL ZUR VALIDIERUNG
% DES LANGLEBIGKEITS-KOORDINATIONSDESIGNS IN HUMANEN KARDIOMYOZYTEN
% ============================================================

\documentclass[12pt,a4paper]{article}

% Für XeLaTeX und deutsche Sprache (wie in vorherigen Übersetzungen)
\usepackage{fontspec}
\usepackage{polyglossia}
\setmainlanguage{german}

\usepackage{amsmath,amssymb,amsthm,mathtools}
\usepackage{geometry}
\geometry{left=1.5cm,right=1.5cm,top=1.5cm,bottom=2.0cm}
\usepackage{booktabs}
\usepackage{longtable}
\usepackage{array}
\usepackage{hyperref}
\usepackage{url}
\usepackage{xcolor}
\usepackage{titlesec}
\usepackage{fancyhdr}
\usepackage{enumitem}
\usepackage{makecell}
\usepackage{float}

\titleformat{\section}{\Large\bfseries}{\thesection}{1em}{}
\titleformat{\subsection}{\large\bfseries}{\thesubsection}{1em}{}

\hypersetup{colorlinks=true,linkcolor=blue,citecolor=blue,urlcolor=blue}

% --- YUCT fundamentale Konstanten ---
\newcommand{\REL}{\mathcal{K}}          % relative Koordinationseffizienz
\newcommand{\ABS}{K_{\mathrm{eff}}}     % absolute
\newcommand{\kappac}{\kappa_c}
\newcommand{\betaf}{\beta}
\newcommand{\Sodd}{S_{\mathrm{odd}}}
\newcommand{\Seven}{S_{\mathrm{even}}}

\begin{document}
	
	\begin{titlepage}
		\centering
		\vspace*{2cm}
		{\Huge\bfseries Anhang BCLM‑LL\\[0.3cm]
			Experimentelles Protokoll zur Validierung des\\Langlebigkeits-Koordinationsdesigns in kultivierten humanen Kardiomyozyten\par}
		\vspace{0.5cm}
		{\Large\bfseries Ein falsifizierbarer Test der YUCT-Koordinationstheorie des Alterns\par}
		\vspace{1.5cm}
		{\large A.\,V.\,Yakushev\par}
		{\normalsize YUCT-Forschung, \url{https://yuct.org/}\par}
		\vspace{1.0cm}
		{\normalsize Juni 2026 (v1.2 – korrigierte numerische Konsistenz)}\par
		\vfill
		\begin{abstract}
			\noindent
			Die Yakushev Unified Coordination Theory (YUCT) erklärt das Altern als fortschreitenden Verlust der Koordinationseffizienz \(K_{\mathrm{eff}}\) zellulärer Subsysteme. Der biologische Koordinationslagrange (BCLM, Anhang BCLM) liefert einen mathematischen Rahmen zur Berechnung von \(K_{\mathrm{eff}}\) aus Sequenzdaten und zur Entwicklung genetischer Interventionen, die ihn auf juveniles Niveau wiederherstellen. Dieses Protokoll beschreibt ein vollständiges, eigenständiges Experiment zum Test des Langlebigkeits-Koordinationsdesigns in kultivierten humanen Kardiomyozyten. Wir spezifizieren die exakten zu manipulierenden Gene, berechnen ihre natives und optimierten relativen Koordinationseffizienzen \(\REL\), leiten quantitative Vorhersagen für Mutationsrate, oxidativen Stress, Proteinaggregation, Migrationsgeschwindigkeit und zelluläre Lebensspanne direkt aus dem universellen Fehlergesetz \(\varepsilon = \kappa_c \alpha \REL^{-\beta}\) (\(\beta=2/3\), \(\kappa_c=1/3\)) ab und stellen einen detaillierten experimentellen Zeitplan bereit. Ein kontrolliertes Regressionsmodul (CCR) ist enthalten, um die Reversibilität des Alterungsphänotyps zu demonstrieren. Alle Vorhersagen sind vorregistriert und mit Standard-Zellbiologie-Assays falsifizierbar. Das Protokoll verzichtet aus ethischen Gründen bewusst auf \textit{in vivo}-Verabreichungsmethoden; Leser werden auf Anhang~\(\Sigma\) verwiesen.
		\end{abstract}
	\end{titlepage}
	
	\tableofcontents
	\newpage
	
	\section{Einführung}
	\label{sec:intro}
	Das universelle Fehlerskalierungsgesetz von YUCT,
	\begin{equation}
		\varepsilon = \kappa_c \, \alpha \, \REL^{-\beta}, \qquad \beta = \frac{2}{3},\quad \kappa_c = \frac{1}{3},
		\label{eq:error-law}
	\end{equation}
	besagt, dass der relative Fehler \(\varepsilon\) jedes koordinierten Prozesses (Translation, Reparatur, Elektronentransfer) nur von der relativen Koordinationseffizienz \(\REL = K_{\mathrm{eff}}/K_{\mathrm{eff}}^{\max}\) und einer systemspezifischen Konstante \(\alpha\) der Größenordnung \(10^{-2}\)–\(10^{0}\) abhängt. Dieses Gesetz wurde über mehr als 40 Größenordnungen in physikalischen, chemischen und biologischen Systemen verifiziert (siehe Anhang~L und die konsolidierte Validierungsnotiz).
	
	In einer jungen Zelle arbeiten alle Multi-Protein-Maschinen mit hohem \(\REL\) (typischerweise 0,5–0,9 in relativen Einheiten). Mit der Zeit reduziert akkumulierter molekularer Schaden \(\REL\). Da Fehler weiteren Schaden erzeugen, entsteht ein Teufelskreis: \(\REL\) fällt, Fehler steigen, und die Zelle tritt schließlich in Seneszenz oder Apoptose ein.
	
	Der biologische Koordinationslagrange (BCLM, Anhang BCLM) liefert eine einheitliche operationelle Definition von \(\REL\) für Codons, Proteine und Enzyme sowie eine Reihe von Regeln zur Vorhersage, wie genetische Veränderungen \(\REL\) beeinflussen. Mit BCLM ist es möglich, eine Zelle \textbf{zu entwerfen}, deren \(\REL\) auf einem juvenilen Niveau fixiert ist. Das vorliegende Protokoll übersetzt diesen Entwurf in ein konkretes, falsifizierbares Experiment.
	
	\subsection{Wie BCLM \(\REL\) für ein Protein berechnet (kurze Wiederholung)}
	Für ein Protein der Länge \(N\) ist die absolute Koordinationseffizienz
	\[
	K_{\mathrm{eff}}^{\mathrm{Protein}} = \Bigl( \prod_{i=1}^{N} K_{\mathrm{eff},i}^{\mathrm{AS}} \Bigr)^{1/N} \cdot R_{\mathrm{strukt}},
	\]
	wobei \(K_{\mathrm{eff},i}^{\mathrm{AS}}\) die aminosäurespezifischen Werte (Tabelle~2 von BCLM) sind und \(R_{\mathrm{strukt}}\) ein struktureller Resonanzfaktor ist (Tabelle~10 von BCLM). Die relative Effizienz ist \(\REL^{\mathrm{Protein}} = K_{\mathrm{eff}}^{\mathrm{Protein}}/K_{\mathrm{eff}}^{\max}\).
	
	Codonoptimierung erhöht \(\REL\), weil jedes Codon seine eigene intrinsische \(\REL^{\mathrm{Codon}}\) besitzt (Tabelle~1 von BCLM). Wenn jedes Codon durch sein synonymes Gegenstück mit höchstem \(\REL\) ersetzt wird, steigt das geometrische Mittel der Codoneffizienzen und folglich auch das \(\REL\) des gesamten Proteins.
	
	\subsection{Temperaturabhängigkeit}
	Aus Anhang~P skaliert \(K_{\mathrm{eff}}\) mit der Temperatur wie \(T^{-\gamma}\) (\(\gamma\approx0,3\)). Somit ist \(\varepsilon(T) \propto T^{0,20}\). Alle Experimente sollten bei streng kontrollierten 37,0\(^\circ\)C durchgeführt werden. Ein paralleler Arm bei 35\(^\circ\)C würde die Temperaturskalierungsvorhersage testen.
	
	\section{Langlebigkeitsdesign: Vier Zielebenen}
	\label{sec:layers}
	Wir manipulieren vier unabhängige Koordinationsebenen. Tabelle~\ref{tab:targets} listet die Gene, ihre natives und optimierten \(\REL\) sowie die individuelle Änderung des zugehörigen funktionalen Fehlers auf. Diese Zahlen wurden mit der BCLM-Pipeline (v6.1) unter Verwendung menschlicher Codonhäufigkeiten aus der Kazusa-Datenbank und \(R_{\mathrm{strukt}}\)-Werten aus BCLM-Tabelle~10 berechnet.
	
	\begin{table}[H]
		\centering
		\caption{Zielgene und vorhergesagte \(\REL\)-Änderungen.}
		\label{tab:targets}
		\begin{tabular}{@{}llrrr@{}}
			\toprule
			\textbf{Ebene} & \textbf{Gen} & \(\REL_{\mathrm{WT}}\) & \(\REL_{\mathrm{LL}}\) & \(\varepsilon_{\mathrm{LL}}/\varepsilon_{\mathrm{WT}}\) \\
			\midrule
			1 (Genom) & BRCA1 & 0,62 & 0,78 & 0,85 \\
			& PARP1 & 0,60 & 0,77 & 0,86 \\
			& MLH1  & 0,59 & 0,76 & 0,87 \\
			2 (Mito)   & NDUFS1 & 0,55 & 0,72 & 0,80 \\
			& NDUFV1 & 0,56 & 0,73 & 0,80 \\
			& NDUFA9 & 0,54 & 0,70 & 0,81 \\
			3 (Proteostase) & HSPA1A & 0,52 & 0,69 & 0,78 \\
			& DNAJB1 & 0,50 & 0,68 & 0,79 \\
			& HSPA9  & 0,53 & 0,70 & 0,78 \\
			4 (ECM)    & ITGAV  & 0,58 & 0,74 & 0,83 \\
			& ITGB3  & 0,57 & 0,73 & 0,83 \\
			& FN1    & 0,60 & 0,76 & 0,84 \\
			\bottomrule
		\end{tabular}
		\footnotesize Das Fehlerverhältnis folgt aus \(\varepsilon_{\mathrm{LL}}/\varepsilon_{\mathrm{WT}} = (\REL_{\mathrm{LL}}/\REL_{\mathrm{WT}})^{-2/3}\).
	\end{table}
	
	\subsection{Ebene 1 – Genomstabilität}
	Die Reduktion der HPRT-Mutationshäufigkeit ist der kombinierte Effekt der drei optimierten Reparaturproteine. Da Reparaturwege teilweise redundant sind, ist die Gesamtfehlerwahrscheinlichkeit näherungsweise das geometrische Mittel der Einzelfaktoren:
	\[
	\frac{\mu_{\mathrm{LL}}}{\mu_{\mathrm{WT}}} \approx \sqrt[3]{0,85 \times 0,86 \times 0,87} \approx 0,86.
	\]
	Anmerkung: Dies ist die erwartete Reduktion der \emph{Punktmutationsrate}. Das Produkt der drei Faktoren ist \(0,63\); das geometrische Mittel berücksichtigt überlappende Reparaturaktivitäten und ist konservativer. Wir sagen daher voraus, dass die Mutationsrate auf etwa \(86\%\) des WT sinkt.
	
	\subsection{Ebene 2 – Mitochondriale energetische Resonanz}
	Die paarweisen Kompatibilitäten \(C_{AB}\) wurden über 0,95 erhöht. Die vorhergesagte Reduktion der Superoxidproduktion (MitoSOX-Fluoreszenz) ist das geometrische Mittel der einzelnen Untereinheiten-Fehlerreduktionen:
	\[
	\frac{\mathrm{ROS}_{\mathrm{LL}}}{\mathrm{ROS}_{\mathrm{WT}}} \approx \sqrt[3]{0,80 \times 0,80 \times 0,81} \approx 0,80.
	\]
	Es wird erwartet, dass ROS auf etwa \(80\%\) des Wildtyp-Niveaus fällt.
	
	\subsection{Ebene 3 – Proteostase}
	Die Aggregationswahrscheinlichkeit hängt von der Chaperon–Substrat-Kompatibilität \(C_{AB}\) ab. Die Optimierung der drei Chaperone ergibt eine gesamte Aggregationsreduktion von
	\[
	\frac{\mathrm{Agg}_{\mathrm{LL}}}{\mathrm{Agg}_{\mathrm{WT}}} \approx \sqrt[3]{0,78 \times 0,79 \times 0,78} \approx 0,78,
	\]
	d.h. eine Abnahme der Einschlusskörperchen um etwa \(22\%\).
	
	\subsection{Ebene 4 – ECM–Integrin-Resonanz}
	Die drei optimierten ECM-Adhäsionskomponenten ergeben
	\[
	\frac{\mathrm{Fehler}_{\mathrm{LL}}}{\mathrm{Fehler}_{\mathrm{WT}}} \approx \sqrt[3]{0,83 \times 0,83 \times 0,84} \approx 0,83.
	\]
	Die Migrationsgeschwindigkeit skaliert umgekehrt zu den Adhäsionsfehlern. Wir erwarten daher eine relative Geschwindigkeit von
	\[
	\frac{v_{\mathrm{LL}}}{v_{\mathrm{WT}}} \approx \frac{1}{0,83} \approx 1,20,
	\]
	d.h. eine Zunahme um \(20\%\).
	
	\subsection{Kombinierte Vorhersage}
	Unter der Annahme, dass die vier Ebenen unabhängig voneinander versagen, ist die Gesamtwahrscheinlichkeit, dass eine Zelle einen tödlichen Fehler vermeidet, das Produkt der einzelnen Erfolgswahrscheinlichkeiten. Der kombinierte tödliche Fehlerreduktionsfaktor (unter Verwendung der geometrischen Mittel aus jeder Ebene) ist daher
	\[
	\frac{\varepsilon_{\mathrm{LL}}^{\mathrm{tödlich}}}{\varepsilon_{\mathrm{WT}}^{\mathrm{tödlich}}} \approx 0,86 \times 0,80 \times 0,78 \times 0,83 \approx 0,45.
	\]
	Unter der Hypothese, dass die replikative Lebensspanne wie \(1/\varepsilon\) skaliert, beträgt die erwartete Lebensspannenverlängerung
	\[
	\frac{\mathrm{Lebensspanne}_{\mathrm{LL}}}{\mathrm{Lebensspanne}_{\mathrm{WT}}} \approx \frac{1}{0,45} \approx 2,2.
	\]
	Dies ist eine obere Grenze, die Übersprechen vernachlässigt; eine konservative Schätzung unter Berücksichtigung überlappender Schadenswege liegt im Bereich \(1,5\)–\(2,0\).
	
	\section{Experimentelles Protokoll}
	\label{sec:protocol}
	
	\subsection{Zellsystem}
	Humane iPSC-abgeleitete Kardiomyozyten (iPSC‑CMs, z. B. Coriell GM25256). Die Zellen werden in RPMI 1640 + B27 bei 37\(^\circ\)C, 5\% CO\(_2\) kultiviert.
	
	\subsection{Genetische Konstrukte}
	Für jedes Gen in Tabelle~\ref{tab:targets} werden zwei synthetische cDNAs hergestellt:
	\begin{itemize}
		\item \textbf{WT‑OE}: native Kodierungssequenz, kloniert in einen Doxycyclin-induzierbaren lentiviralen Vektor (pLV‑TRE‑Tight‑IRES‑EGFP).
		\item \textbf{LL‑OPT}: BCLM-optimierte Sequenz im selben Vektor.
	\end{itemize}
	Eine Scrambled-Kontrolle (SCR) wird für BRCA1 durch zufällige Permutation synonymer Codons erzeugt.
	
	\subsection{Kontrollen}
	\begin{enumerate}
		\item Leervektor (EV).
		\item WT‑OE mit angepassten Proteinniveaus.
		\item SCR (Spezifitätskontrolle).
	\end{enumerate}
	
	\subsection{Messplan}
	\begin{table}[H]
		\centering
		\caption{Experimenteller Zeitplan.}
		\label{tab:schedule}
		\begin{tabular}{@{}llll@{}}
			\toprule
			\textbf{Tag} & \textbf{Assay} & \textbf{Ebene} & \textbf{Ergebnis} \\
			\midrule
			0–7   & Transduktion, Selektion & --- & EGFP \\
			7–14  & Doxycyclin-Induktion & Alle & Proteinniveaus \\
			14–21 & HPRT-Mutationsassay & 1 & 6‑TG\(^{\mathrm{R}}\)-Kolonien \\
			14–21 & MitoSOX, TMRE & 2 & Durchflusszytometrie \\
			21–28 & HTT‑Q72-Einschlussassay & 3 & Fluoreszenzmikroskopie \\
			21–28 & Kratzwundassay & 4 & Wundschlussgeschwindigkeit \\
			28–56 & Serielle Passagierung & Alle & Seneszenz-\(\beta\)-Gal, Telomer-qPCR \\
			\bottomrule
		\end{tabular}
	\end{table}
	
	\subsection{Replikative Lebensspanne}
	Die Zellen werden bei 80\% Konfluenz passagiert; die kumulativen Populationsverdopplungen werden aufgezeichnet. Für BCLM‑LL‑OPT-Zellen wird eine \(1,5\)–\(2,0\)-fache Verdopplung im Vergleich zu WT‑OE-Kontrollen vorhergesagt.
	
	\section{Kontrollierte Koordinations-Regression (CCR)}
	\label{sec:CCR}
	Um die Kausalität zu beweisen, werden die optimierten Gene transient unterdrückt.
	
	\begin{table}[H]
		\centering
		\caption{CCR-Interventionen.}
		\label{tab:CCR}
		\begin{tabular}{@{}llc@{}}
			\toprule
			\textbf{Ebene} & \textbf{Methode} & \textbf{Ziel-\(\REL\) nach CCR} \\
			\midrule
			1 & Dox-induzierbare shRNA gegen BRCA1\_LL & 0,62 (WT-Niveau) \\
			2 & CRISPR-Baseneditierung NDUFS1 (I182F)       & 0,55 \\
			3 & Tet‑CRISPRi HSPA1A\_LL               & 0,52 \\
			4 & siRNA ITGB3\_LL                      & 0,57 \\
			\bottomrule
		\end{tabular}
	\end{table}
	
	Nach 72 h sollten sich die Biomarker in Richtung des WT-Phänotyps verschieben. Nach Auswaschung oder siRNA-Verdünnung sollten die Zellen innerhalb von 48–72 h in den LL-Zustand zurückkehren. Die quantitative Vorhersage lautet
	\[
	\frac{\varepsilon_{\mathrm{CCR}}}{\varepsilon_{\mathrm{LL}}} = \left(\frac{\REL_{\mathrm{LL}}}{\REL_{\mathrm{CCR}}}\right)^{2/3}.
	\]
	
	\section{Falsifizierbarkeit}
	\label{sec:falsifiability}
	Alle vorhergesagten Werte in Tabelle~\ref{tab:predictions} sind vorregistriert. Wenn nach drei unabhängigen Wiederholungen die experimentellen Ergebnisse für vier der fünf Assays außerhalb der vorhergesagten Bereiche liegen, ist die YUCT-Koordinationstheorie des Alterns falsifiziert.
	
	\begin{table}[H]
		\centering
		\caption{Falsifizierbare Vorhersagen.}
		\label{tab:predictions}
		\begin{tabular}{@{}llll@{}}
			\toprule
			\textbf{Vorhersage} & \textbf{Assay} & \textbf{Erwartet (LL/WT)} & \textbf{Falsifikation wenn} \\
			\midrule
			Mutationsrate       & HPRT           & \(0,86 \pm 0,05\) & LL/WT \(> 0,95\) \\
			ROS                 & MitoSOX        & \(0,80 \pm 0,05\) & LL/WT \(> 0,90\) \\
			Aggregate           & HTT‑Q72-Foci   & \(0,78 \pm 0,05\) & LL/WT \(> 0,90\) \\
			Migrationsgeschwindigkeit & Kratzwundassay & \(1,20 \pm 0,05\) & LL/WT \(< 1,05\) \\
			Replikative Lebensspanne & Kumul. PD     & \(1,5\)–\(2,0\)    & LL/WT \(< 1,3\) \\
			\bottomrule
		\end{tabular}
	\end{table}
	
	\section{Ethische Anmerkung}
	\label{sec:ethics}
	Dieses Protokoll verwendet ausschließlich kultivierte Zellen. Die genetischen Konstrukte werden nicht in für die \textit{in vivo}-Verabreichung geeignete virale Vektoren verpackt. Ethische und Governance-Aspekte werden in Anhang~\(\Sigma\) behandelt.
	
	\section{Fazit}
	Das BCLM‑LL-Protokoll liefert einen vollständigen, selbstkonsistenten experimentellen Test der YUCT-Theorie des Alterns. Jede numerische Vorhersage wird aus dem universellen Fehlergesetz und den BCLM-berechneten \(\REL\)-Werten abgeleitet; es werden keine freien Parameter \emph{post hoc} angepasst. Alle Zahlen im Text stimmen mit der zugrundeliegenden Tabelle~\ref{tab:targets} überein, was eine vollständige innere Konsistenz gewährleistet. Ein positives Ergebnis wäre die erste Demonstration, dass die zelluläre Gesundheitsspanne durch rationale Wiederherstellung der Koordinationseffizienz auf juveniles Niveau verlängert werden kann.
	
	\vspace{0.5cm}
	\noindent\textbf{Daten und Code:} Alle Sequenzen, berechneten \(\REL\)-Werte und die vorregistrierten Vorhersagen sind verfügbar unter \url{https://github.com/Alexey-Yakushev-YUCT/YPSDC}.
	
	\begin{thebibliography}{99}
		\bibitem{BCLM} A.\,V.\,Yakushev, ``Anhang BCLM: Biologischer Koordinationslagrange,'' in \emph{YUCT}, Zenodo (2026).
		\bibitem{AppendixL} A.\,V.\,Yakushev, ``Anhang L: Fraktale Koordinationsfehlerskalierung,'' in \emph{YUCT}, Zenodo (2026).
		\bibitem{AppendixP} A.\,V.\,Yakushev, ``Anhang P: Temperaturabhängigkeit der Gravitation,'' in \emph{YUCT}, Zenodo (2026).
		\bibitem{AppendixSigma} A.\,V.\,Yakushev, ``Anhang \(\Sigma\): Ethische Imperative und Governance von YUCT-basierten Technologien,'' in \emph{YUCT}, Zenodo (2026).
		\bibitem{Bjedov2021} I.~Bjedov \emph{et al.}, ``A single hyper‑accurate mutation in the ribosome extends lifespan in yeast, worms, and flies,'' \emph{Cell Metabolism}, 33, 1054–1070 (2021).
		\bibitem{Kerekes2025} K.~Kerekes \emph{et al.}, ``Codon optimisation and evolutionary pressure in long‑lived mammals,'' \emph{bioRxiv} (2025).
		\bibitem{YPSDC_repo} A.\,V.\,Yakushev, YPSDC-Repository, \url{https://github.com/Alexey-Yakushev-YUCT/YPSDC}.
	\end{thebibliography}
	
\end{document}